% ============================================================
% Lab W1 -- SUBMISSION TEMPLATE
%
% This is a template, not the assignment itself. It mirrors the
% structure of Lab W1 (labw1-intro-setup.tex) so you can fill in
% your group's answers directly here and submit this file (plus
% the PDF it produces).
%
% HOW TO USE THIS FILE:
%   1. Keep this file in the same folder as coa-accessible.sty
%      (the same folder the Lab W1 PDF and .tex file came in).
%   2. Replace anything that says [Type your answer here.] with
%      your actual answer, then delete the italics if you want.
%   3. Compile with pdflatex TWICE in a row. The first run writes
%      the page-count information, the second run reads it back in
%      so "Page X of Y" is correct. If you are using an editor like
%      Overleaf or TeXstudio, it usually does this for you
%      automatically when you click "build" or "recompile."
%
% Lines starting with a percent sign, like this one, are comments.
% LaTeX ignores everything after a % on a line, so comments never
% show up in your PDF. Feel free to delete these tip comments once
% you don't need them anymore, or just leave them; they are harmless.
% ============================================================
\documentclass[11pt]{article}
\usepackage{coa-accessible}

% This line sets the title that appears at the top of the PDF and
% in the document's metadata. The first {...} is the main title,
% the second {...} is the smaller line underneath it.
\AssignmentMeta{Lab W1 Submission: Environment Setup}{Group submission template}

\begin{document}
\AssignmentTitleBlock
% Note: this template intentionally skips the "AI note" and
% "Resources note" boxes that appear on the assignment sheet.
% Those describe how the assignment itself was written; this file
% is your own work, so they do not apply here.

\section*{Group members}
% LaTeX TIP: \begin{itemize} ... \end{itemize} makes a bulleted
% list. Each \item is one bullet point. Replace the bracketed text
% below with your actual names.
%
% One LaTeX quirk worth knowing: if \item is immediately followed
% by a square bracket, like \item [text], LaTeX assumes you are
% giving that item a custom label instead of a bullet, and the
% bullet disappears. That is why the brackets below are written as
% {[} and {]} instead of plain [ and ]; wrapping them in braces
% tells LaTeX "this is just a character," not "here comes an
% option." You will only need this trick when a line happens to
% start with a square bracket right after \item.
\begin{itemize}
  \item {[}Type name here.{]}
  \item {[}Type name here.{]}
  \item {[}Type name here.{]}
  \item {[}Type name here.{]}
\end{itemize}

\section{Part A -- QEMU checklist}
% LaTeX TIP: this table environment was already built for you in
% coa-accessible.sty, so you do not need to know how tables work
% yet to use it. Just edit the text between the & symbols. Each &
% moves to the next column, and \\ ends a row. Type Yes or No (or
% a short note) in the second column for each row.
\begin{accessibletable}{Part A checklist}{tab:parta-submission}{p{0.65\linewidth}p{0.25\linewidth}}
\toprule
Checklist item & Confirmed by all 4? \\
\midrule
QEMU installed and \texttt{qemu-system-x86\_64 --version} runs & [Yes/No] \\
AArch64 cross-compiler installed & [Yes/No] \\
Native hello-world compiles and runs & [Yes/No] \\
Cross-compiled AArch64 hello-world runs under QEMU & [Yes/No] \\
\bottomrule
\end{accessibletable}

\section{Part B -- Raspberry Pi output}
Paste your \texttt{uname -m} output below.
% LaTeX TIP: the lstlisting environment below displays text in a
% monospace font with a box around it, exactly as typed, which is
% perfect for pasting terminal output. Do not put % comment
% characters inside it; everything between \begin and \end is
% shown exactly as-is.
\begin{lstlisting}
[Paste your uname -m output here.]
\end{lstlisting}

Paste the relevant line(s) from \texttt{cat /proc/cpuinfo} below (the
\texttt{model name} or \texttt{Hardware} field is enough, you do not need
the whole file).
\begin{lstlisting}
[Paste your /proc/cpuinfo output here.]
\end{lstlisting}

\section{Screenshots}
% LaTeX TIP: \accessiblefigure{file}{width}{caption}{label} inserts
% an image. The steps below are commented out (each line starts
% with %) so this template compiles cleanly even before you add
% your screenshots. To use them:
%   1. Save your screenshots as .png or .jpg files in this same
%      folder, and give them short names with no spaces, for
%      example qemu-setup.png and pi-setup.png.
%   2. Delete the % at the start of the \accessiblefigure lines
%      below.
%   3. Change the file name in the first {} to match your actual
%      file, and edit the caption text (the third {}) to describe
%      what the screenshot shows. That caption doubles as
%      alt text, so make it genuinely descriptive.

% \accessiblefigure{qemu-setup.png}{0.8\linewidth}{Terminal output showing the cross-compiled AArch64 hello-world program running successfully under QEMU.}{fig:qemu}

% \accessiblefigure{pi-setup.png}{0.8\linewidth}{Terminal output showing uname -m and the hello-world program running natively on the Raspberry Pi.}{fig:pi}

[If you have not added your screenshot files yet, leave a short note
here about what you still need to attach.]

\section{Part C -- Reflection}
% LaTeX TIP: \begin{enumerate} ... \end{enumerate} makes a numbered
% list instead of a bulleted one. The question text is left in
% place below so it stays clear which answer is which; type your
% answer right after each question.
\begin{enumerate}
  \item What did \texttt{uname -m} report on the Raspberry Pi, and what does
        that string actually name?

        \textit{[Type your answer here.]}
  \item You cross-compiled for AArch64 on your laptop in Part A but compiled
        natively on the Pi in Part B. In your own words, why does a
        cross-compiler exist at all? What problem does it solve that a
        regular compiler cannot?

        \textit{[Type your answer here.]}
  \item QEMU let you run an AArch64 binary on an x86-64 laptop. Based on what
        you observed, what do you think QEMU is actually doing to make that
        possible?

        \textit{[Type your answer here.]}
\end{enumerate}

\section*{Anything that did not work}
If any step did not work for someone in the group, describe what you tried
here. If everything worked, you can leave this section as-is or delete it.

\textit{[Type your answer here, or delete this section.]}

\end{document}
